\input{preamble}

\section*{Crossing symmetry}

We have for Bhabha scattering
\begin{equation*}
\langle|\mathcal{M}|^2\rangle=2\alpha^2
\left(\frac{s^2+u^2}{t^2}+\frac{2u^2}{st}+\frac{t^2+u^2}{s^2}\right)
\end{equation*}

and for Moller scattering
\begin{equation*}
\langle|\mathcal{M}|^2\rangle=2\alpha^2
\left(\frac{s^2+u^2}{t^2}+\frac{2s^2}{tu}+\frac{s^2+t^2}{u^2}\right)
\end{equation*}

Bhabha and Moller scattering share $s\leftrightarrow u$ crossing symmetry.

\bigskip
We have for electron-positron annihilation
\begin{equation*}
\langle|\mathcal{M}|^2\rangle=2\alpha^2
\left(\frac{u}{t}+\frac{t}{u}\right)
\end{equation*}

and for Compton scattering
\begin{equation*}
\langle|\mathcal{M}|^2\rangle=-2\alpha^2
\left(\frac{u}{s}+\frac{s}{u}\right)
\end{equation*}

Annihilation and Compton scattering share $s\leftrightarrow t$ crossing symmetry.
The minus sign is the result of moving a fermion from inbound (in annihilation)
to outbound (in Compton).

\bigskip
\href{https://georgeweigt.github.io/examples/crossing-symmetry-demo.html}{Eigenmath script}

\end{document}
